
\documentclass[12pt, a4paper]{article}
\usepackage[utf8]{inputenc}
\usepackage[T2A]{fontenc}
\usepackage{amssymb}
\usepackage[left=2cm, right=2cm, top=2cm, bottom=2cm]{geometry}

\begin{document}

Вопрос №5. Что такое множество нулевой меры на прямой, а на плоскости? Привести примеры. Показать, что множество алгебраических чисел бесконечно, но имеет нулевую меру. Ответ. Говорят, что множество меры $0$, лежащее в $\mathbb{R}$, если его можно накрыть набором интервалов. Примеры на прямой: целые числа $\mathbb{Z}$, рациональные числа $\mathbb{Q}$, множество из одной точки $x_0$. 

Пример один. Прямая. Пусть $A = \{a_n\} \subset \mathbb{R}$ — множество корней полинома $P(x) = a_n x^n + \dots + a_1 x + a_0$, где $a_i$ — целые числа. Каждому полиному соответствует набор $\{a_0, a_{-1}, ..., a_{n-1}\}$, а количество корней не более $n$. Объединение конечного числа счётных множеств счётно. 

Таким образом, множество алгебраических чисел, как объединение счётного числа счётных множеств, имеет нулевую меру.
\end{document}
